\subsection{DERIVATIVE WITH RESPECT TO A PARAMETER OF AN INTEGRAL OVER A COMPACT INTERVAL}

Let A be a compact neighbourhood of a point \(\alpha_{0}\) in the field R (resp. the field C), let I = {[}a, b{]} be a compact interval in R, and f a continuous map of I × A into a complete normed space E over R (resp. C). We have seen (II, p.~69, prop. 2) that under these conditions \(\mathbf{g}(\alpha) = \int_{a}^{b} \mathbf{f}(t, \alpha) dt\) is a continuous function on A. Let us seek sufficient conditions for g to admit a derivative at the point \(\alpha_{0}\) . One has, for \(\alpha \neq \alpha_{0}\) ,

\[
\frac {\mathbf {g} (\alpha) - \mathbf {g} (\alpha_ {0})}{\alpha - \alpha_ {0}} = \int_ {a} ^ {b} \frac {\mathbf {f} (t , \alpha) - \mathbf {f} (t , \alpha_ {0})}{\alpha - \alpha_ {0}} d t
\]

so (II, p.~69, cor. 2), if the functions \(x \mapsto \frac{\mathbf{f}(x, \alpha) - \mathbf{f}(x, \alpha_0)}{\alpha - \alpha_0}\) converge uniformly on I to a (necessarily continuous) function \(x \mapsto \mathbf{h}(x)\) as \(\alpha\) tends to \(\alpha_0\) (while remaining \(\neq \alpha_0\) ), then \(\mathbf{g}\) admits a derivative equal to \(\int_{a}^{b} \mathbf{h}(t) dt\) at the point \(\alpha_0\) ; moreover, for each \(x \in I\) , \(\frac{\mathbf{f}(x, \alpha) - \mathbf{f}(x, \alpha_0)}{\alpha - \alpha_0}\) tends to \(\mathbf{h}(x)\) , so \(\mathbf{h}(x)\) is the derivative at the point \(\alpha_0\) of the map \(\alpha \mapsto \mathbf{f}(x, \alpha)\) ; we denote this derivative (called the partial derivative of \(\mathbf{f}\) with respect to \(\alpha\) ) by the notation \(\mathbf{f}_{\alpha}'(x, \alpha_0)\) ; the hypotheses we have made imply that

\[
\mathbf {g} ^ {\prime} (\alpha_ {0}) = \int_ {a} ^ {b} \mathbf {f} _ {\alpha} ^ {\prime} (t, \alpha_ {0}) d t.\tag{12}
\]

The following proposition gives a very simple sufficient condition for the validity of formula (12):

PROPOSITION 5. Suppose that the partial derivative \(\mathbf{f}_{\alpha}^{\prime}(x,\alpha)\) exists for all \(x\in \mathrm{I}\) and all \(\alpha\) in an open neighbourhood \(\mathbf{V}\) of \(\alpha_0\) , and that, for all \(\alpha \in \mathbf{V}\) , the map \(x\mapsto \mathbf{f}_{\alpha}^{\prime}(x,\alpha)\) is regulated on \(\mathrm{I}\) . Under these conditions, if \(x\mapsto \mathbf{f}_{\alpha}^{\prime}(x,\alpha)\) converges uniformly on \(\mathrm{I}\) to \(x\mapsto \mathbf{f}_{\alpha}^{\prime}(x,\alpha_0)\) as \(\alpha\) tends to \(\alpha_0\) , then the function \(\mathbf{g}(\alpha) = \int_{a}^{b}\mathbf{f}(t,\alpha)dt\) admits a derivative, given by the formula (12), at the point \(\alpha_0\) .

Indeed, for every \(\varepsilon > 0\) there exists by hypothesis an \(r > 0\) such that \(|\alpha - \alpha_0| \leqslant r\) implies \(\left\| \mathbf{f}_{\alpha}'(x, \alpha) - \mathbf{f}_{\alpha}'(x, \alpha_0) \right\| \leqslant \varepsilon\) for any \(x \in I\) . By props. 3 and 5 of I, p.~17 one has, for \(|\alpha - \alpha_0| \leqslant r\) ( \(\alpha \neq \alpha_0\) ) and for all \(x \in I\)

\[
\left| \left| \frac {\mathbf {f} (x , \alpha) - \mathbf {f} (x , \alpha_ {0})}{\alpha - \alpha_ {0}} - \mathbf {f} _ {\alpha} ^ {\prime} (x, \alpha_ {0}) \right| \right| \leqslant \varepsilon
\]

which proves the uniform convergence of \(\frac{\mathbf{f}(x,\alpha)-\mathbf{f}(x,\alpha_{0})}{\alpha-\alpha_{0}}\) to \(\mathbf{f}_{\alpha}^{\prime}(x,\alpha_{0})\) on I as \(\alpha\) tends to \(\alpha_{0}\) (remaining \(\neq\alpha_{0}\) ), and so establishes formula (12).

COROLLARY. If the partial derivative \(\mathbf{f}_{\alpha}^{\prime}(x,\alpha)\) exists on I × V and is a continuous function of \((x,\alpha)\) on this set, then the function g admits a derivative given by the formula (12) at the point \(\alpha_{0}\) .

Indeed, if W is a compact neighbourhood of \(\alpha_{0}\) contained in V, then the map \((x,\alpha)\mapsto\mathbf{f}_{\alpha}^{\prime}(x,\alpha)\) is uniformly continuous on the compact set I × W, so \(\mathbf{f}_{\alpha}^{\prime}(x,\alpha)\) tends to \(\mathbf{f}_{\alpha}^{\prime}(x,\alpha_{0})\) uniformly on I as \(\alpha\) tends to \(\alpha_{0}\) .

From prop. 5 one deduces a more general proposition which allows one to evaluate the derivative of an integral when, not only the integrand f, but also the limits of integration, depend on a parameter \(\alpha\) :

PROPOSITION 6. Supposing that the hypotheses of prop. 5 are satisfied, let \(a(\alpha)\) , \(b(\alpha)\) be two functions defined on V, with values in I; if the derivatives \(a'(\alpha_0)\) , \(b'(\alpha_0)\) exist and are finite then the function \(\mathbf{g}(\alpha) = \int_{a(\alpha)}^{b(\alpha)} \mathbf{f}(t, \alpha) dt\) admits at \(\alpha_0\) a derivative given by the formula

\[
\mathbf {g} ^ {\prime} (\alpha_ {0}) = \int_ {a (\alpha_ {0})} ^ {b (\alpha_ {0})} \mathbf {f} _ {\alpha} ^ {\prime} (t, \alpha_ {0}) d t + b ^ {\prime} (\alpha_ {0}) \mathbf {f} (b (\alpha_ {0}), \alpha_ {0}) - a ^ {\prime} (\alpha_ {0}) \mathbf {f} (a (\alpha_ {0}), \alpha_ {0}).\tag{13}
\]

Indeed, for all \(\alpha \in \mathbf{V}\) distinct from \(\alpha_0\) one can write

\[
\begin{array}{c} \frac {\mathbf {g} (\alpha) - \mathbf {g} (\alpha_ {0})}{\alpha - \alpha_ {0}} = \int_ {a (\alpha_ {0})} ^ {b (\alpha_ {0})} \frac {\mathbf {f} (t , \alpha) - \mathbf {f} (t , \alpha_ {0})}{\alpha - \alpha_ {0}} d t + \frac {1}{\alpha - \alpha_ {0}} \int_ {b (\alpha_ {0})} ^ {b (\alpha)} \mathbf {f} (t, \alpha) d t \\ - \frac {1}{\alpha - \alpha_ {0}} \int_ {a (\alpha_ {0})} ^ {a (\alpha)} \mathbf {f} (t, \alpha) d t. \end{array}
\]

By prop. 5 of II, p.~74, the first integral on the right-hand side tends to \(\int_{a(\alpha_0)}^{b(\alpha_0)} \mathbf{f}'_\alpha(t, \alpha_0) dt\) as \(\alpha\) tends to \(\alpha_0\) . In the second integral we replace \(\mathbf{f}(t, \alpha)\) by \(\mathbf{f}(b(\alpha_0), \alpha_0)\) and show that the difference tends to 0. We put \(\mathbf{M} = \operatorname{Max}\big( \| \mathbf{f}(b(\alpha_0), \alpha_0) \|, |b'(\alpha_0)| + 1 \big)\) ; the function \(b(\alpha)\) being continuous at the point \(\alpha_0\) and the function \(\mathbf{f}\) continuous at the point \((b(\alpha_0), \alpha_0)\) , for every \(\varepsilon\) such that \(0 < \varepsilon < 1\) there exists an \(r > 0\) such that the relation \(|\alpha - \alpha_0| \leqslant r\) entails that \(\| \mathbf{f}(t, \alpha) - \mathbf{f}(b(\alpha_0), \alpha_0) \| \leqslant \varepsilon\) for all \(t\) belonging to the interval with endpoints \(b(\alpha_0)\) and \(b(\alpha)\) ; thus one may also suppose that the relation \(|\alpha - \alpha_0| \leqslant r\) entails \(\left| \frac{b(\alpha) - b(\alpha_0)}{\alpha - \alpha_0} - b'(\alpha_0) \right| \leqslant \varepsilon\) .

By the mean value formula (II, p.~62, formula (17)) one thus has

\[
\left\| \frac {1}{\alpha - \alpha_ {0}} \int_ {b (\alpha_ {0})} ^ {b (\alpha)} \mathbf {f} (t, \alpha) d t - \frac {b (\alpha) - b (\alpha_ {0})}{\alpha - \alpha_ {0}} \mathbf {f} (b (\alpha_ {0}), \alpha_ {0}) \right\| \leqslant \left| \frac {b (\alpha) - b (\alpha_ {0})}{\alpha - \alpha_ {0}} \right| \varepsilon
\]

and consequently

\[
\left\| \frac {1}{\alpha - \alpha_ {0}} \int_ {b (\alpha_ {0})} ^ {b (\alpha)} \mathbf {f} (t, \alpha) d t - b ^ {\prime} (\alpha_ {0}) \mathbf {f} (b (\alpha_ {0}), \alpha_ {0}) \right\| \leqslant 2 \mathrm{M} \varepsilon
\]

which shows that \(\frac{1}{\alpha - \alpha_0} \int_{b(\alpha_0)}^{b(\alpha)} \mathbf{f}(t, \alpha) dt\) tends to \(b'(\alpha_0)\mathbf{f}(b(\alpha_0), \alpha_0)\) . In the same way one shows that \(\frac{1}{\alpha - \alpha_0} \int_{a(\alpha_0)}^{a(\alpha)} \mathbf{f}(t, \alpha) dt\) tends to \(a'(\alpha_0)\mathbf{f}(a(\alpha_0), \alpha_0)\) .

